\section*{Chapter \ref{ch:g7:prop} --- Proportionality}

\begin{solution}{exo:g7:prop:1}
First table: \omterm{def:g3:division:remainder}{quotients} $\frac{7.5}{3} = 2.5$,
$\frac{17.5}{7} = 2.5$, $\frac{27.5}{11} = 2.5$: \omterm{def:g6:propdata:def}{proportional},
coefficient $2.5$.

Second table: $\frac62 = 3$ and $\frac{15}{5} = 3$, but
$\frac{28}{9} \approx 3.1$: not \omterm{def:g6:propdata:def}{proportional}.
\end{solution}

\begin{solution}{exo:g7:prop:2}
With the cross rule: $5 \times p = 6 \times 8$, so
$p = \frac{48}{5} = 9.60$ euros.
\end{solution}

\begin{solution}{exo:g7:prop:3}
Rate: $36 \div 3 = 12$ pages per minute. In $5$ minutes:
$60$ pages. For $96$ pages: $96 \div 12 = 8$ minutes.
\end{solution}

\begin{solution}{exo:g7:prop:4}
$30\,\%$ of $250$: $0.3 \times 250 = 75$. \quad
$15\,\%$ of $60$: $0.15 \times 60 = 9$. \quad
$t\,\%$ of $80$ is $20$: $\frac{20}{80} = \frac{25}{100}$, so
$t = 25$.
\end{solution}

\begin{solution}{exo:g7:prop:5}
$\frac{18}{25} = \frac{18 \times 4}{25 \times 4} = \frac{72}{100} =
72\,\%$.
\end{solution}

\begin{solution}{exo:g7:prop:6}
Path: $9 \times 25\,000 = 225\,000$ cm $= 2.25$ km.
Lake: $2$ km $= 200\,000$ cm on the ground, so
$200\,000 \div 25\,000 = 8$ cm on the map.
\end{solution}

\begin{solution}{exo:g7:prop:7}
$4.3$ m $= 430$ cm, and $430 \div 43 = 10$ cm.
\end{solution}

\begin{solution}{exo:g7:prop:8}
Chocolate per serving: $240 \div 8 = 30$ g. With $300$ g:
$300 \div 30 = 10$ servings exactly.
\end{solution}

\begin{solution}{exo:g7:prop:9}
\emph{1.} Discount: $20$ euros, i.e.\ $\frac{20}{80} = 25\,\%$ of the
original price.

\emph{2.} The rise of $20$ euros is now compared with the \emph{new}
reference $60$: $\frac{20}{60} = \frac13 \approx 33\,\%$. A percentage
always refers to a whole; the whole changed, so the percentage does
too.
\end{solution}

\begin{solution}{exo:g7:prop:10}
Coefficient from the complete column: $\frac{17.5}{10} = 1.75$. Then
$y = 4 \times 1.75 = 7$ and $x = \frac{10.5}{1.75} = 6$.
\end{solution}

\begin{solution}{exo:g7:prop:11}
\emph{1.} After one copy: $10 \times 0.8 = 8$ cm. After two:
$8 \times 0.8 = 6.4$ cm.

\emph{2.} The second reduction applies to the already-reduced copy,
not to the original: the factors multiply,
$0.8 \times 0.8 = 0.64$, so two copies reduce to $64\,\%$ of the
original --- not $60\,\%$.
\end{solution}

\begin{solution}{pb:g7:prop:1}
\textbf{1.} Heights and shadows are \omterm{def:g6:propdata:def}{proportional}, so by the
cross rule (\cref{thm:g7:prop:cross}), with the stick's column
$(1,\ 0.4)$ and the tree's $(h,\ 3.2)$:
$0.4 \times h = 1 \times 3.2$, hence
$h = \frac{3.2}{0.4} = 8$ m.

\textbf{2.} Tower: $\frac{26}{0.4} = 65$ m. Child's shadow:
$1.5 \times 0.4 = 0.6$ m (shadow $=$ height $\times$ the
moment's coefficient $0.4$).

\textbf{3.} As the sun moves across the sky, the coefficient
linking heights to shadows changes; only measurements taken at
the same moment share the same coefficient.

\textbf{4.} At that moment the coefficient is $1$: every height
equals its shadow. The \omterm{def:g5:solids:def}{pyramid}'s apex is therefore
$147$ m high --- the height of the Great \omterm{def:g5:solids:def}{Pyramid} (its shadow's
tip, measured from the point below the apex, is all one needs).

\textbf{5.} The column \omterm{def:g3:division:remainder}{quotient}
$\frac{\text{shadow}}{\text{height}}$ --- the \omterm{def:g7:prop:table}{proportionality}
coefficient of the moment --- is common to every vertical
object: that is exactly what makes the table of heights and
shadows a \omterm{def:g7:prop:table}{proportionality table} (\cref{def:g7:prop:table}).

\textbf{6.} Sun rays arrive \omterm{def:g6:lines:perp}{parallel}. On a \emph{flat} Earth,
two \omterm{def:g6:lines:perp}{parallel} rays would strike two vertical sticks at the same
angle: both would cast \omterm{def:g6:propdata:def}{proportional} shadows --- either both no
shadow, or both a shadow. One stick with no shadow (Syene) and
one with a shadow (Alexandria), at the same instant, is
impossible on a flat Earth: the two verticals must point in
different directions --- the surface curves.

\textbf{7.} On the drawing, the Syene vertical points straight
at the sun; the Alexandria vertical is tilted by the angle
between the two city directions, seen from the center. Since the
rays are \omterm{def:g6:lines:perp}{parallel}, that tilt is exactly the $7.2^\circ$ measured
between ray and stick in Alexandria.

\textbf{8.} $\frac{7.2}{360} = \frac{72}{3600} = \frac{1}{50}$:
one fiftieth of a full turn.

\textbf{9.} Arc lengths are \omterm{def:g6:propdata:def}{proportional} to angles: if
$7.2^\circ$ --- one fiftieth of the turn --- corresponds to
$800$ km, the full turn corresponds to
\[
50 \times 800 = 40\,000 \text{ km}.
\]

\textbf{10.} \omterm{def:g4:geometry:circle}{Diameter} $= \frac{\text{circumference}}{\pi}
\approx \frac{40\,000}{3.14} \approx 12\,700$ km --- against
the modern $12\,742$ km: correct to well within one percent,
with a stick.

\textbf{11.} $40\,000$ km $= 4\,000\,000\,000$ cm, and
$4\,000\,000\,000 \div 40\,000\,000 = 100$ cm $= 1$ m of
circumference. \omterm{def:g4:geometry:circle}{Diameter}: $100 \div 3.14 \approx 32$ cm --- a
handsome desk globe.

\textbf{12.} $4.8$ km $= 480\,000$ cm, and
$480\,000 \div 40\,000\,000 = 0.012$ cm $= 0.12$ mm. The
highest mountain of the Alps is a tenth of a millimeter on a
\omterm{def:g4:numbers:round}{meter-round} globe: a speck of dust. At this \omterm{def:g7:prop:scale}{scale} the Earth is
smoother than most polished balls.

\textbf{13.} $40\,000 \div 40 = 1\,000$ days, and
$1\,000 \div 365 \approx 2.7$: nearly three years of walking.

\textbf{14.} $\frac{10}{6\,370} \approx 0.0016$, that is about
$0.2\,\%$ (more precisely $0.16\,\%$): the atmosphere is a
whisper-thin skin on the planet.

\textbf{15.} Error: $42\,000 - 40\,000 = 2\,000$ km, and
$\frac{2\,000}{40\,000} = \frac{5}{100} = 5\,\%$. One sentence:
with a stick, a well, a measured road and one \omterm{def:g7:prop:table}{proportionality},
Eratosthenes measured a planet to within a few percent ---
mathematics travels far on very little.
\end{solution}
